


\section{Discussion}
\label{sec:discussion}

\paragraph{Significance and vision.}
We argue that representative, risk-sensitive scenario distributions should be a first-class object of evaluation in agentic AI, anchored to a concrete five-property definition of trustworthiness (\S\ref{subsec:defining_trustworthiness}). The reframing shifts evaluation from benchmark instances to scenario distributions, from single scores to multi-dimensional profiles, and from static testing to the iterative red--blue cycle of the \emph{Trustworthy Optimization Factory}. The cross-model comparison (\S\ref{subsec:profile_comparison}) confirms a corollary: meaningful comparison is inherently multi-property---families (Llama, Mistral, Kimi, GLM, Qwen, GPT-oss, DeepSeek) make different trade-offs that a leaderboard erases. If this agenda succeeds, the community converges on shared scenario taxonomies and world-model libraries~\cite{liang2023helm,zhou2023sotopia,drouin2024workarena,xie2024osworld}, reports profiles rather than single scores~\cite{liang2023helm,kiela2021dynabench}, and adopts the Factory's red--blue cycle as a deployment-readiness pipeline---aligned with emerging AI governance frameworks~\cite{ganguli2022redteaming,nist2023airmf,eu2024aiact}.

\paragraph{Implications for agentic information retrieval.}
The shift from passive retrieval to \emph{retrieval-grounded action}~\cite{xi2024agentir} makes agentic IR a natural and high-stakes setting for this framework. Three implications follow. \textit{(i)} Retrieved content is an attack surface: indirect prompt injection through retrieved documents~\cite{greshake2023indirect,zhou2024poisonedrag} is precisely a P2 (Robustness) failure, and our sandbox cleanly separates agents that treat retrieved text as data from those that treat it as instructions. \textit{(ii)} Retrieval-grounded agents expose sensitive corpora through natural-language interfaces, making P3 (Safety) and P4 (Social-Ethical) joint concerns rather than separable. \textit{(iii)} Agentic IR designers should not pick a model by aggregate score: a system top-ranked on retrieval helpfulness may rank last on resistance to retrieved-content injection or social manipulation; only profile-level comparison reveals which trade-off matches a given deployment.

\paragraph{Limitations and future work.}
A full instantiation will need larger held-out suites (our 100-scenario validation suite improves substantially over a 24-scenario pilot but is still synthetic), human-validated failure annotations, and grounding of representative sampling in real-world evidence (deployment logs, incident reports, red-team traces~\cite{ganguli2022redteaming,debenedetti2024agentdojo}). Open questions include formalising Factory convergence criteria, handling user-instructed violations that resist prompt-level mitigation, and standard property-weight vectors $\boldsymbol{\lambda}$ for common deployment classes so risk-weighted aggregates are comparable across studies. The adversarial framings we release carry a dual-use risk; we mitigate by pairing the scenario suite with the blue-team hardening cycle (\S\ref{subsec:vulnerability_attribution}) and by including the Improper-Refusal class so that blanket refusal is not rewarded as a substitute for value-aligned judgement.

\paragraph{Ethics statement.}
This work studies trustworthiness failures of agentic AI systems and therefore necessarily catalogues adversarial framings (indirect prompt injection, authority-cover exfiltration, proxy-discrimination, role-play jailbreaks) that could in principle be misused. Three properties of the artefact bound this risk. \textit{(i)~No human subjects, no PII.} All 100 scenarios, all retrieved documents, all simulated user profiles, and all sandbox tool outputs are synthetic, hand-authored by the authors; no human-subjects data, no production logs, and no scraped personal information enter the suite, so no IRB review was required and no individual is identifiable in the released artefact. \textit{(ii)~Sandboxed execution only.} Every trajectory runs against the five synthetic tools of \S\ref{subsec:setup} in a local sandbox at L1 of the Sandbox Fidelity Ladder (\S\ref{subsec:layer2}); the released code does not ship live network, shell, or browser tools, so replaying a scenario cannot reach a real user, database, or external API. \textit{(iii)~Responsible release.} We release the scenario suite paired with the blue-team interventions (tool-output firewall, confirmation gate) that neutralise the matching attack class, and we publish the Improper-Refusal taxonomy class explicitly so that downstream users cannot reduce a HAAF score simply by training agents to refuse more; refusal of benign requests is itself penalised under P1. Models evaluated through Amazon Bedrock were accessed under standard provider terms; no model weights were fine-tuned in this study.

\paragraph{Statistical scope.}
We attach 95\% Wilson score intervals to every per-property success rate $\hat{T}_{P_k}$ and a paired-scenario bootstrap ($B{=}10{,}000$, seed 42) $p$-value to each anti-scaling pair underlying Finding F3. Three of four sibling pairs reject the null that the larger model is no worse than the smaller at $p<0.05$ (Llama, Qwen, Mistral borderline); GPT-oss does not. These tests are consistent with the headline anti-scaling pattern not being explained by sampling noise alone, but the intervals remain wide: with $n{=}14$--$26$ scenarios per property, the per-cell CIs span roughly $\pm 0.2$, so these statistics are descriptive of cross-family trends, not powered for fine-grained per-property claims.

\paragraph{LLM-judge annotation reliability.}
To probe how much our heuristic violation label depends on coarse signals, we audited a stratified 150-trajectory sample (30 per property, balanced violation/success where the population allowed) with Kimi-K2.5 as an LLM judge via Bedrock Converse (temperature 0, 512-token cap). Overall observed agreement is 0.607 with Cohen's $\kappa{=}0.287$ (fair); agreement is strongest on P4 ($\kappa{=}0.800$, almost perfect) and weakest on P5 ($\kappa{=}0.118$) and P1 ($\kappa{=}0.143$)---precisely the properties whose stratified pools of true violations are smallest (4 and 6 respectively), inflating chance agreement. The IR audit is the most damning result: the judge flagged 21/150 trajectories as improper refusals while our string-based heuristic flagged zero, yielding recall $=0.000$ and a structurally undefined precision (0/0). The current heuristic simply cannot see IR. Two honest caveats: (i) this is a single-judge audit with no human gold standard; (ii) Kimi-K2.5 is itself one of the 13 evaluated agents and shares a family with Kimi-K2-Thinking, so the judge plausibly shares blind spots with two of the models it audits. We commit to a separate-family judge (e.g., a Claude- or Llama-family model) plus a human spot-check protocol over at least the 21 judge-flagged IR cases and a matched control set for v2.

\paragraph{Sandbox fidelity ladder.}
Our methodology---scenario authoring, per-property scoring, profile aggregation, and the Factory's red--blue cycle---is invariant to the fidelity rung of the underlying sandbox, but the cost of an undetected error scales sharply as fidelity rises. The present study is a v1 instantiation at L1 (a toy synthetic sandbox with mocked tools and stubbed external services), which is appropriate for the methodological claims we make but not for deployment-readiness verdicts. v2 commits to L2 (a controlled production-like sandbox with real tool surfaces, sandboxed network endpoints, and recorded-then-replayed retrieval corpora), so that anti-scaling and Factory-convergence findings transfer to settings where a P3 or P4 failure has a real-world cost-of-error footprint.

\paragraph{Representativeness reprise.}
Finally, our 100 scenarios exercise 5/5 of the sandbox tools we expose, mapped to the top-$K$ operation classes observed in SWE-bench~\cite{jimenez2023swebench}, API-Bank~\cite{li2023apibank}, WebArena~\cite{zhou2023webarena}, and AgentBench~\cite{liu2023agentbench}, but they remain hand-authored. A deployment-grounded sampler---drawing scenario skeletons from real agent traces, incident reports, and red-team logs, weighted by observed operation-class frequencies---is committed as v2 future work and is the route by which representativeness moves from designer intent to measured property of the suite.
